\documentclass[12pt]{article}

\usepackage{amsmath}
\usepackage{amssymb}
\usepackage{amsthm}
\usepackage[margin=1in]{geometry}
\usepackage{hyperref}
\usepackage{booktabs}
\usepackage{url}

\hypersetup{
  colorlinks=true,
  linkcolor=blue,
  urlcolor=blue,
  citecolor=blue
}

\newtheorem{principle}{Principle}

\begin{document}

\begin{titlepage}
\centering
\vspace*{2cm}

{\LARGE\bfseries The End of Armchair Ethics\\[0.5em]
\large\textit{How Invariance Testing Makes Normative Systems Falsifiable}}

\vspace{1.5cm}

{\Large Andrew H. Bond, Senior Member, IEEE}

\vspace{0.5cm}

{\large
Department of Computer Engineering\\
San Jos\'e State University\\
San Jos\'e, CA 95192, USA\\[0.3cm]
\texttt{andrew.bond@sjsu.edu}\\[0.3cm]
ORCID: 0009-0003-2599-6158
}

\vspace{1cm}

{\large Spring 2026}

\vspace{2cm}

\begin{abstract}
For twenty-five centuries, ethical claims have been unfalsifiable. You cannot run an
experiment to determine whether utilitarianism is correct. No observation settles
whether Kant was right. This paper argues that the question was malformed. The
answerable question is not whether an ethical theory is true, but whether an ethical
judgment system is consistent, non-gameable, and accountable. These are empirical
properties with pass/fail criteria. We present a framework---Philosophy
Engineering---that provides the first falsifiability conditions for normative systems:
declared invariances, operational tests, minimal witnesses for failures, and machine-checkable
audit artifacts. The result is not a new ethical theory. It is the
infrastructure that makes ethical engineering subject to the same discipline as any
other engineering domain.
\end{abstract}

\vfill

\noindent\rule{\linewidth}{0.4pt}
\small
\noindent\textbf{AI-Assisted Writing Disclosure:} Portions of this manuscript were drafted and refined with AI assistance (Claude, Anthropic). All technical content, mathematical results, and claims were conceived, verified, and approved by the author.

\noindent\textbf{Ethics Statement:} This research involves no human subjects. Empirical validation uses publicly available corpora and synthetic data. No IRB approval was required.

\noindent\textbf{Data Availability:} Code and data are available at \url{https://github.com/ahb-sjsu/erisml-lib}.

\noindent\textbf{Conflict of Interest:} The author declares no conflicts of interest. This research received no external funding.

\noindent\textbf{CRediT Statement:} Andrew H. Bond: Conceptualization, Methodology, Formal Analysis, Investigation, Writing --- Original Draft, Writing --- Review \& Editing, Software.
\noindent\rule{\linewidth}{0.4pt}
\normalsize

\end{titlepage}

\noindent\textbf{Keywords:} falsifiability, ethics, invariance testing, normative systems,
philosophy engineering, gauge theory, consistency, accountability, AI alignment

\section{The Twenty-Five Century Stalemate}

Consider the state of ethical discourse. Utilitarians argue with deontologists. Virtue
ethicists critique both. Contractarians propose alternatives. The debate continues,
century after century, because there is no empirical arbiter. No experiment can be
run. No observation is decisive. The arguments are sophisticated, the intuitions are
marshaled, but the stalemate is permanent.

This is not a failure of ethics. It is a category confusion. We have been asking a
question that may have no answer---``Which ethical theory is true?''---while ignoring
questions that demonstrably do.

Compare: ``Is Euclidean geometry true?'' is not quite the right question. The better
questions are: Is Euclidean geometry internally consistent? Does it apply to this
domain? Do its predictions match observation in the relevant regime? These are
answerable. The first is provable. The second and third are empirical.

Ethics has lacked this separation. We conflated the (possibly unanswerable) question
of ethical truth with the (demonstrably answerable) questions of ethical system
integrity. This paper concerns only the latter.

\section{The Testable Questions}

An ethical judgment system---whether implemented in a human institution or an
AI---takes representations of situations as input and produces evaluations as output.
About such systems, we can ask:

\begin{description}
  \item[Consistency:] Does the system produce the same judgment for the same
    situation, regardless of how that situation is described?
  \item[Non-gameability:] Can the system be exploited by redescribing situations in
    semantically equivalent but superficially different ways?
  \item[Accountability:] When judgments differ, can the system attribute the
    difference to a change in the situation, a change in the evaluative
    commitments, or uncertainty?
  \item[Non-triviality:] Does the system actually distinguish between genuinely
    different situations, or does it collapse everything to a constant output?
\end{description}

These are not metaphysical questions. They are engineering questions. They have
operational definitions. They produce yes/no answers. They generate witnesses when
the answer is no.

\section{The Core Principle}

The central construct is the Epistemic Invariance Principle (EIP):

\begin{principle}[Epistemic Invariance Principle]
A judgment procedure is epistemically well-posed only if it is
invariant---up to declared output equivalence---under a declared set of
meaning-preserving transformations.
\end{principle}

Unpacked: if you declare that certain transformations of input (renaming variables,
permuting option order, paraphrasing, changing units) should not change the
judgment, then the system must actually be invariant under those transformations. If
it is not, you have a witness: a specific transformation that flips the output. This
witness is a reproducible, minimal counterexample. It is evidence of a defect.

This is falsifiability. Not for ethical truth---that may remain forever beyond reach---but
for ethical system integrity. And system integrity is what we can engineer.

\section{The Geometry of Consistency}

There is a deeper structure here. The requirement that evaluation be invariant under
a group of meaning-preserving transformations is precisely the structure studied in
gauge theory---the mathematics of symmetry and invariance that underlies modern
physics. This geometric perspective is developed fully in~\cite{bond2026ge}.

This is not metaphor. The formal apparatus is identical:

\begin{quotation}
Representations live in a space. Transformations act on that space. Invariance
means the evaluation function factors through the quotient---it depends only
on the equivalence class, not the representative.
\end{quotation}

\begin{description}
  \item[Curvature] measures the failure of invariance to be path-independent. Non-zero
    curvature means there exist sequences of transformations that return to
    the ``same'' representation but produce different evaluations. This is an
    exploit. It is detectable.
  \item[Conservation laws] (when applicable) provide monitored quantities that
    should remain stable. Drift signals symmetry-breaking or model mismatch.
\end{description}

The physics vocabulary is not essential. What is essential is the operational content:
invariance requirements can be tested, violations can be witnessed, and the
geometric structure provides diagnostic tools for finding exploitable inconsistencies.

\section{What This Is Not}

This framework does not:

\begin{description}
  \item[Tell you which values to hold.] The choice of what matters is a governance
    problem, not an engineering problem. The framework tests whether your
    chosen commitments are applied consistently, not whether they are correct.
  \item[Prove ethical truths.] We remain agnostic about whether there are ethical
    truths in any metaphysical sense. The framework operates at the level of
    systems, not reality.
  \item[Eliminate all failure modes.] Sensor spoofing, implementation bugs,
    incomplete grounding, adversarial inputs outside the declared envelope---these
    remain. The framework localizes where risk lives; it does not eliminate
    risk.
  \item[Replace human judgment.] It provides discipline for judgment systems,
    human or artificial. The discipline is: declare your invariances, test them,
    produce witnesses when they fail, audit everything.
\end{description}

\section{The Paradigm Shift}

The shift can be summarized in one table:

\begin{table}[h]
\centering
\begin{tabular}{ll}
\toprule
\textbf{Traditional Ethics} & \textbf{Philosophy Engineering} \\
\midrule
Is this action wrong? & Does this judgment flip under meaning-preserving transforms? \\
Is utilitarianism correct? & Is this utilitarian system internally consistent and non-gameable? \\
Argue from intuitions & Produce minimal witnesses for failures \\
Debate endlessly & Run the test suite \\
Status: unfalsifiable & \textbf{Status: pass/fail} \\
\bottomrule
\end{tabular}
\end{table}

The right column does not answer the questions in the left column. It replaces them
with answerable questions. This is not a retreat. It is a clarification of what can be
known.

\section{The Pragmatist Foundation}

This framework rests on a pragmatist epistemology: formal systems---logic,
mathematics, modal semantics---are tools, not mirrors of metaphysical structure.
They are judged by their utility, coherence, and predictive power, not by their alleged
correspondence to a mind-independent reality.

This is not skepticism. It is epistemic discipline. The null hypothesis is: do not posit
metaphysical structures until evidence forces you. Evidence would be evidence.
Logical arguments from modal systems are not evidence. Intuitions about necessity
are not evidence. The arguments may be interesting, but they do not meet the
burden.

Under this view, the gauge-theoretic structure of invariance requirements is not a
discovery that ``ethics has shape'' in some metaphysical sense. It is the observation
that if you want representation-invariance, gauge-theoretic tools are the appropriate
formalism---just as arithmetic is appropriate if you want to count. The tool fits the
problem. That is all.

\section{Implementation}

The framework is implementable. The key components are:

\begin{description}
  \item[Transformation registries:] Explicit, versioned, hashed declarations of
    which input transformations are meaning-preserving.
  \item[Test suites:] Generators that produce transformed inputs spanning the
    declared invariances.
  \item[Canonicalizers:] Functions that map inputs to normal forms within
    equivalence classes, enforcing invariance by construction.
  \item[Witness reducers:] Algorithms that minimize failing transformations to
    produce compact, reproducible counterexamples.
  \item[Audit artifacts:] Machine-checkable records of what invariances were
    declared, what tests were run, what witnesses were produced, and what lens
    (evaluative commitments) was applied.
\end{description}

These are not aspirational. They are specified in sufficient detail to be built. The
accompanying technical whitepaper provides JSON schemas, theorem statements,
and worked examples. A reference implementation is now available as ErisML/DEME v3.0 with tensorial ethics support~\cite{bond2026erisml}.

\section{The Forest and the Trees}

It is easy to miss what is happening here. The formalism looks like mathematics. The
vocabulary borrows from physics. The implementation details are software
engineering. Each piece, in isolation, is familiar.

The whole is not familiar. For twenty-five centuries, normative claims have been
debated without any discipline of falsification. Ethical systems have been proposed,
critiqued, refined, and abandoned based on argumentative force, intuitive
plausibility, and rhetorical skill. There was no test suite. There were no witnesses.
There was no audit.

This framework provides those things. Not for ethical truth---that may be beyond
us---but for ethical systems. And in an era of autonomous AI systems that make
normative judgments at scale, the integrity of those systems is what matters
practically.

The question is no longer ``Is this AI's ethics correct?'' That question may be
unanswerable. The question is: ``Does this AI's judgment system flip under
semantically equivalent redescriptions? Can it be gamed? Can its failures be
witnessed and audited?'' These questions have answers. We now have the
infrastructure to find them.

\section{Conclusion}

Philosophy has often been accused of making no progress. In ethics, the accusation
has force: the debates of antiquity are recognizably continuous with the debates of
today. This is not because ethicists lack intelligence or rigor. It is because the
questions they ask---questions about ethical truth---may not admit of empirical
resolution.

The contribution of this work is to separate the unanswerable from the answerable.
We do not solve the problem of ethical truth. We dissolve the conflation that made
ethical system integrity seem equally intractable. It is not.

Ethical systems can be tested for consistency, non-gameability, accountability, and
non-triviality. These tests produce pass/fail verdicts. Failures produce witnesses.
Witnesses enable debugging. Debugging enables improvement. Improvement is
progress.

This is what it looks like when philosophy becomes engineering.

\begin{thebibliography}{9}

\bibitem{bond2025electro}
A.~H.~Bond, ``The Electrodynamics of Value: Gauge-Theoretic Structure in AI
Alignment,'' Technical report, San Jos\'{e} State University, 2025.

\bibitem{bond2025phileng}
A.~H.~Bond, ``Philosophy Engineering: A Technical Whitepaper on the Epistemic
Invariance Principle,'' Technical report, San Jos\'{e} State University, 2025.

\bibitem{bond2025pragmatist}
A.~H.~Bond, ``A Pragmatist Rebuttal to Logical and Metaphysical Arguments for God,''
Technical report, San Jos\'{e} State University, 2025.

\bibitem{bond2026ge}
A.~H.~Bond, \emph{Geometric Ethics: The Mathematical Structure of Moral Reasoning}, v1.14, Department of Computer Engineering, San Jos\'e State University, Spring 2026.

\bibitem{bond2026erisml}
A.~H.~Bond, ``ErisML/DEME v3.0: A Library for Governed, Foundation-Model-Enabled Agents with Tensorial Ethics,'' \url{https://github.com/ahb-sjsu/erisml-lib}, 2026.

\end{thebibliography}

\end{document}
